\newpage
\section{Introduction}
% ──────────────────────────────────────────────────────────────────────────────

Explainable Artificial Intelligence (XAI) encompasses a set of methods designed
to render machine learning models interpretable by revealing internal
mechanisms and the reasoning processes underlying their
decisions~\cite{Arrieta2020XAI}. In clinical medicine, interpretability is not
merely desirable but increasingly mandated: clinical decision-making carries
direct consequences for patient safety, and regulatory frameworks across
multiple jurisdictions now require that AI-assisted diagnoses be accompanied by
interpretable justifications~\cite{jin2023guidelines}. Despite this regulatory
momentum, the metrics used to evaluate explanation quality in medical AI remain
poorly standardised. Current evaluation methodologies rely heavily on
subjective user studies which, while valuable for assessing clinician
perception, are susceptible to confirmation bias and frequently fail to assess
the technical fidelity of the generated
explanations~\cite{mirzaei2023xai}.

This standardisation deficit is compounded by the prevalence of post-hoc
explanation methods in clinical AI. Because such explanations are generated
after model training, they may not faithfully reflect the model's actual
decision-making logic. In medical contexts, this gap is particularly
consequential: a persuasive but unfaithful explanation could reinforce a
clinician's incorrect judgement rather than correct it, potentially
compromising patient outcomes.

The present study focuses on the medical domain and addresses three
research questions: (RQ1)~Which quantitative evaluation metrics are
currently employed to assess XAI explanation quality in medical AI?
(RQ2)~How do these metrics co-occur across studies, and which pairs
exhibit potential redundancy? (RQ3)~Where do systematic coverage gaps
persist across medical sub-domains? By analysing a focused corpus of 52
peer-reviewed medical XAI studies spanning five sub-domains---oncology,
cardiology, neuroimaging, radiology, and clinical decision
support---we quantify metric co-occurrence patterns, identify
redundancies in current practice, and propose a minimal but comprehensive
evaluation framework grounded in a three-layer taxonomy of mathematical
faithfulness, system robustness, and human-centred trust.

The principal contributions of this paper are threefold:
\begin{enumerate}
    \item A systematic cross-study co-occurrence analysis quantifying the distribution, pairing patterns, and correlations of XAI evaluation metrics across five medical sub-domains.
    \item A three-layer evaluation taxonomy---mathematical faithfulness, system robustness, and human-centred trust---grounded in empirical usage patterns and validated against the 52-study corpus.
    \item A minimal evaluation framework specifying one metric per evaluative dimension, designed to ensure comprehensive coverage without imposing excessive methodological burden.
\end{enumerate}


\section{Background}
% ──────────────────────────────────────────────────────────────────────────────

Artificial intelligence is being rapidly integrated into clinical workflows
across diagnostic imaging, prognostic modelling, and therapeutic decision
support. However, because complex AI models frequently operate as opaque
systems, clinicians cannot readily ascertain how diagnostic or prognostic
outputs are generated. This opacity creates a pressing need for XAI methods
that render medical AI systems transparent and
interpretable~\cite{Tjoa2021MedicalXAI}. XAI has been applied to domains
including clinical decision
support~\cite{abbas2025explainable}, breast cancer
diagnosis~\cite{saharan2025explainable}, cardiovascular risk
prediction~\cite{shah2025explainable}, and responsible cardiovascular disease
screening~\cite{muneer2025blockchain}.

A central distinction in XAI is between models that are interpretable by
design (intrinsically interpretable models) and methods that generate
explanations after training (post-hoc methods). Post-hoc approaches are
appealing because they can be applied to virtually any deployed clinical
system without requiring architectural modification. However, this
flexibility entails a fundamental limitation: there is no guarantee that the
explanation accurately represents the model's internal computations, and
this fidelity gap remains one of the core unresolved challenges in medical
XAI~\cite{holzinger2022information,palatnik2021xai}.

Two broad paradigms have emerged for evaluating explanation quality in
medicine. \emph{Technical evaluation} employs quantitative metrics---such as
faithfulness, sensitivity, robustness, and consistency---to assess how
accurately an explanation approximates actual model behaviour.
\emph{Human-centred evaluation}, by contrast, focuses on the end user,
measuring constructs such as clinician trust, cognitive load, and practical
usefulness through observational studies and behavioural
experiments~\cite{lenis2020domain}.

While both paradigms have matured considerably, they remain largely
disconnected in medical practice. A method that achieves high scores on
technical metrics may produce explanations that are opaque or misleading to
clinicians, while an explanation that non-specialists find intuitive may
lack mathematical fidelity to the underlying model. This disconnect
impedes the establishment of shared evaluation standards and limits the
comparability of results across medical sub-domains---a problem this
paper directly addresses through cross-domain co-occurrence analysis
and structured taxonomy construction.

% ──────────────────────────────────────────────────────────────────────────────
